\cleardoubleoddpage
\chapter{Experimental Setup}
\label{chap:experiments}

\section{Evidence Tiers and Frozen Scope}

The repository separates engineering evidence from thesis evidence. Synthetic smoke runs test
that the data contract, encoder, cache, sequence trainer, rollout evaluator, and report writers
execute without private data or a GPU. Small real-data checks validate the Cyclone/KvikIO schema
and target fields. Only the byte-verified medium universe and its frozen protocol enter the
scientific comparison. This separation is important: a successful command is not evidence that a
surrogate generalizes.

The accepted comparison is a retrospective five-fold nested group cross-validation study. The
51-trajectory universe was available for development and its manifest is therefore not a pristine
locked test set. Each outer fold holds out trajectories that are disjoint from its training and
inner-validation groups. Model and checkpoint choices are made inside the fold, before the outer
test evidence for that fold is opened. The reported uncertainty is an estimate of performance on
this finite universe, not a claim of deployment-level generalization.

\section{Medium Dataset and Data Binding}

The consumed universe contains 51 trajectories and 1,173 retained timestep bundles. Every
trajectory contributes 23 retained time steps under the preprocessing contract: four input
channels, spatial shape $(32,8,16,85,32)$, offset 80, temporal subsampling by eight, stored
$\code{kyspec}$ and $\code{fluxspec}$ targets, and a preferred float32 representation. The
channel-first snapshot contract is
\begin{equation}
  x \in \mathbb{R}^{B\times C\times S_1\times S_2\times S_3\times S_4\times S_5}.
\end{equation}
The byte-level universe manifest records the ordered trajectory identifiers, the 1,173 sampled
binary timestep files, and the metadata files required by the loader. Both the durable consumed
subset and the operational public-data alias were checked against the same manifest by hashing
each consumed file; no path is used as a dataset identity.

The five outer folds contain 30 or 31 training trajectories, 10 or 11 validation trajectories,
and 10 or 11 test trajectories. The exact fold manifests are committed under
\path{experiment_protocols/} and are addressed by their SHA-256 digest. Within each fold, all
normalization statistics, encoder parameters, latent caches, and sequence checkpoints are derived
from the fold's training trajectories. This is an end-to-end split: a trajectory used by an
encoder is never called an outer-test trajectory for that fold.

\section{Representation and Temporal Models}

The convolutional encoder maps each preprocessed snapshot to a 128-dimensional latent state. Its
training objective combines a SimSiam-style consistency term with supervised heads for scalar flux
and the two one-dimensional spectra. The heads are retained after encoder training so that a
latent rollout can be evaluated in the original diagnostic space. The work deliberately does not
claim five-dimensional field reconstruction; the decoder required for that claim is outside the
scope of the study.

The sequence candidates are a gated recurrent unit (GRU) and a causal Transformer. Latent-state
persistence is the primary latent-dynamics reference, and direct persistence of the observed flux
is the primary diagnostic baseline. An additional diagnostic-head oracle applies the frozen head
to true future latents. It is an analysis control: it is not a forecast and cannot be interpreted
as an achievable model result. All learned candidates use the same eight-step context, recursive
eight-step forecast horizon, diagnostic lineage, and trajectory-balanced metrics.

Each sequence trainer runs for 500 updates. The GRU has one 256-unit layer; the Transformer uses
width 256, depth three, four attention heads, and no dropout. The cache-normalized variants fit
latent normalization on training trajectories only. Sequence checkpoints are selected by the
lowest validation trajectory-balanced latent RMSE. This checkpoint rule is distinct from family
selection: after each candidate has a validation checkpoint, the GRU or Transformer family with
the lower arithmetic mean of validation trajectory-balanced flux RMSE across seeds is selected
within that outer fold.

\section{Matched Seeds and Stage Ledger}

The development/data seed is fixed at 52. Training seeds are matched across learned candidates at
52, 53, 54, 55, and 56, so initialization variability is visible without changing the trajectory
universe. Persistence controls are evaluated on the same trajectories and horizons and do not
  receive a fictitious training seed. The frozen stage plan contains 255 slots: 230 accepted
  ledger slots (225 metric stages and five selection barriers) and 25 explicitly skipped
  unselected-family test stages. No stage is silently
dropped, retried into the accepted set, or selected from an outer-test metric.

The outer-fold family selections are Transformer for folds 0, 1, and 3, and GRU for folds 2 and
4. This is a fold-wise selection result, not evidence that one architecture is a global winner.
The release manifest records all 25 fold--seed rows, the selected family, stage configuration
hashes, and the exact accepted/skipped counts without publishing trajectory identifiers or server
paths.

\section{Metrics and Estimands}

Flux RMSE is the primary metric because the research question concerns diagnostic forecasting. It
is computed per trajectory and then averaged, so a trajectory with many overlapping windows cannot
dominate a trajectory with fewer valid windows. Latent MSE, flux MAE, spectral relative $L_2$ error,
spectral shape correlation, and finite-output stability are secondary diagnostics. Values are in
preprocessed target units; no physical-unit conversion is claimed.

The primary estimand is the selected learned model minus observed-flux persistence difference in
per-trajectory flux RMSE, averaged hierarchically with equal outer-fold, training-seed, and
trajectory weights. A paired hierarchical bootstrap \citep{saravanan2020hierarchical} resamples training seeds and trajectories
within the fold structure; its 10,000 replicates use a fixed seed recorded in the aggregate. The
secondary estimand is the analogous learned-minus-latent-persistence difference in latent MSE.
We report both point estimates and intervals, plus per-seed variability and a hierarchically
weighted fraction of improved trajectory runs. If the primary interval includes zero, the protocol
requires the conclusion ``no convincing advantage''; a negative result is valid evidence, not a
reason to change the selection rule.

Spectra are reported separately for $\code{kyspec}$ and $\code{fluxspec}$, with relative error and
shape correlation shown together. Their scales and norms differ, so a single pooled spectral
number is not used to claim physical fidelity. No transfer-learning, full-field reconstruction,
long-horizon, solver-speed, or GyroSwin comparison is part of the evidence set.

\section{Reproducibility and W\&B Policy}

The protocol, fold manifests, resolved configuration hashes, stage evidence, and release manifest
form a chain from source tag to result. Raw data, checkpoints, latent caches, and per-trajectory
rows remain outside Git. W\&B is disabled for all 225 accepted metric stages; local status files explicitly
record \code{enabled=false}, \code{requested=false}, and \code{mode=disabled}. Consequently the
thesis cites no live W\&B run or URL. The local manifest, not an online dashboard, is the source of
truth for this result.
